\documentclass[11pt]{article}

\usepackage[T1]{fontenc}
\usepackage{amsmath,amssymb,amsthm}
\usepackage{array}
\usepackage[hidelinks]{hyperref}

\newcommand{\Perm}{\operatorname{Perm}}
\newcommand{\moved}{\operatorname{moved}}
\newcommand{\Equivariant}{\operatorname{Equivariant}}
\newcommand{\Discrete}{\operatorname{Discrete}}
\newcommand{\Finset}{\operatorname{Finset}}
\newcommand{\id}{\operatorname{id}}
\newcommand{\act}{\mathbin{\cdot}}
\newcommand{\swap}[2]{(#1\;#2)}

\newtheorem{proposition}{Proposition}[section]
\theoremstyle{definition}
\newtheorem{definition}[proposition]{Definition}
\newtheorem{example}[proposition]{Example}
\theoremstyle{remark}
\newtheorem{remark}[proposition]{Remark}

\title{A Nominal Package for Lean:\\Permutation, Action and Support Foundations}
\author{}
\date{5 October 2026}

\begin{document}
\maketitle

\begin{abstract}
We present a Lean formalization of
finite permutations, coherent actions, equivariant ordinary functions,
and finite support for elements of selected actions.
Controlled swap factorization keeps every endpoint in the original moved
set and yields factorization avoiding an arbitrary pointwise fixed set.
For any selected action, invariance under swaps outside a set is equivalent
to invariance under all finite permutations fixing that set pointwise.
The design uses Mathlib's group and action interfaces, separates atom and
carrier universes, and makes trivial actions on external data explicit.
Finite bounds are transported by conjugation, preserved by equivariant
images and combined for products. Over infinite atoms, two finite support
bounds have a supporting intersection and each finitely supported element
has a unique least finite support. A Boolean counterexample explains why
finite supportedness alone does not give this conclusion on finite atoms.
Least support commutes with renaming and can shrink under equivariant maps;
a proof-only nominality certificate supplies it uniformly across a carrier.
Counterexamples distinguish actions on the same underlying carrier and show
why action selection and infinitude assumptions are essential to nominal
reasoning. The mathematical interface separates an element's finite dependence
on atoms from the representation of its carrier.
\end{abstract}

\section{Purpose and scope}
\label{sec:introduction}

The intended package supports ordinary Lean proofs about languages with
binders. A researcher should be able to specify syntax, define operations, and
reason about judgments using certified renaming and freshness principles.
The mathematical setting is classical Lean, with one atom sort, actual
quotients, and eventual least supports. Single binders and nested binders form
the initial binding scope. Ordinary propositions remain the language of
theorems; the package is not intended to require a separate object language
for logic.

Before syntax generation or binder reasoning can be certified, the meaning of
renaming must be fixed. A carrier alone does not determine an action: the same
type can carry a trivial action, a canonical atom action, or, for suitable
carriers, a pointwise or conjugation action. These choices affect every later
support statement. The foundation described here therefore exposes a small
permutation and action interface for ordinary Lean proofs.
The underlying mathematics is standard group-action and nominal-set theory;
Pitts~\cite[Chapters 1--2]{pitts2013} provides the background.

For each definition or structure that encodes a mathematical concept, we
state that concept before describing its Lean representation. We identify
the correspondence between the mathematical object, its operations and
laws, and the actual declarations. We also check the pinned Mathlib for
existing constructions before adding package-specific machinery, recording
what is reused and why a different interface is needed when reuse does not
fit. Representation helpers without a separate mathematical counterpart
are described as such, rather than promoted to additional theory.

Declarations use the namespace \texttt{NominalPackage}, under
\path{Package/}, with \path{Package.lean} as
the public import. It uses Mathlib directly and preserves the earlier
\path{Nominal/} and \path{Instances/} development as a reference implementation.
Analogous results in that development are not evidence that a new package
declaration has been proved. The design adapts the subgroup construction for
finite permutations and uses standard \texttt{MulAction} interfaces for the
new boundary.

\paragraph{Implementation status.}
This version describes the implemented foundation increment for roadmap tasks
\texttt{PKG-F02} and the first slice, \texttt{F03a}, of \texttt{PKG-F03}.
The permutation and action interfaces exported by \texttt{Package} have been
checked by Lean, and the source-coverage check and production axiom audit pass.
Section~\ref{sec:evidence} records the verification evidence separately from
the mathematical exposition.

The support-facing slice \texttt{F03b} remains open under \texttt{PKG-F03}.
Supported predicates (\texttt{PKG-01}) are also later work. Persistent usage
examples are reserved for the future case studies.

The present scope includes finite permutations, swaps, actions on atoms,
products and finite atom sets, explicit discrete data, and equivariance of
ordinary functions. It does not establish finite-swap factorization, finite
or least support, freshness, quotient actions, supported predicates, name
abstraction, induction, recursion, or generated syntax. These are later
mathematical layers, with their own proof obligations.

\section{Finite permutations and actions}
\label{sec:foundations}

Let $A$ be the atom carrier. In this section $A$ is an arbitrary type:
neither infinitude nor countability is required. Later existence of fresh
atoms will require infinitude. In Lean, atoms and other carriers have
independent universe parameters,
\[
  A : \texttt{Type}\ u,\qquad X : \texttt{Type}\ v,
  \qquad Y : \texttt{Type}\ w.
\]
The action interface concerns type-valued carriers; no general interface for
\texttt{Sort}-valued carriers is promised.

\subsection{The finite-moved subgroup}
\label{sec:permutations}

\begin{definition}
For a bijection $\pi:A\to A$, define its \emph{moved-point set} by
\[
  \moved(\pi)=\{a\in A\mid \pi(a)\ne a\}.
\]
Let $\Perm(A)$ consist of those bijections whose moved-point sets are finite.
Multiplication is composition in the convention
$(\pi\sigma)(a)=\pi(\sigma(a))$.
\end{definition}

The Lean declaration \texttt{permSubgroup A} is the subgroup of
\texttt{Equiv.Perm A} cut out by this finiteness property, and
\texttt{Perm A} is its carrier. Both lie in the namespace
\texttt{NominalPackage}. The property is a proof field; applying
a permutation does not select a finite enumeration or a least support.
In particular, $\moved(\pi)$ is a set, rather than a chosen
\texttt{Finset}. Equality is extensional: if $\pi(a)=\sigma(a)$ for every
$a$, then $\pi=\sigma$.
The implementation adapts the elementary subgroup argument from
\path{Nominal/Core/FinitePerm.lean}, without importing that module. The
reference revision is recorded in Section~\ref{sec:evidence}.

\begin{proposition}[Subgroup and moved-point laws]
\label{prop:subgroup}
Finite-moved permutations form a subgroup of the bijections of $A$. For
its elements,
\begin{align}
  \moved(1)&=\varnothing,\label{eq:moved-one}\\
  \moved(\pi\sigma)&\subseteq\moved(\pi)\cup\moved(\sigma),
    \label{eq:moved-mul}\\
  \moved(\pi^{-1})&=\moved(\pi),\label{eq:moved-inverse}\\
  \moved(\pi\sigma\pi^{-1})&=\pi[\moved(\sigma)].
    \label{eq:moved-conj}
\end{align}
Here $\pi[S]=\{\pi(a)\mid a\in S\}$ is direct image.
\end{proposition}

\begin{proof}
The identity moves no point. If both $\pi$ and $\sigma$ fix $a$, then so
does $\pi\sigma$, proving~\eqref{eq:moved-mul}. Moreover,
$\pi^{-1}(a)=a$ if and only if $\pi(a)=a$, proving~\eqref{eq:moved-inverse}.
These facts give closure under identity, multiplication, and inverse,
because subsets of finite sets and finite unions are finite. Finally,
injectivity of $\pi$ gives
\[
  (\pi\sigma\pi^{-1})(a)\ne a
  \quad\Longleftrightarrow\quad
  \sigma(\pi^{-1}(a))\ne\pi^{-1}(a).
\]
This is exactly membership of $a$ in the image on the right
of~\eqref{eq:moved-conj}.
\end{proof}

The public application equations are
\[
  1(a)=a,\quad (\pi\sigma)(a)=\pi(\sigma(a)),\quad
  \pi^{-1}(\pi(a))=a,\quad\pi(\pi^{-1}(a))=a.
\]
These are \path{Perm.one_apply}, \path{Perm.mul_apply},
\path{Perm.inv_apply_apply}, and \path{Perm.apply_inv_apply}.
The conversion \path{Perm.toEquiv} preserves identity, multiplication, and
inverse by \path{Perm.toEquiv_one}, \path{Perm.toEquiv_mul}, and
\path{Perm.toEquiv_inv}. Its function application agrees with that of
\texttt{Perm A} by \path{Perm.toEquiv_apply}. A single \texttt{FunLike}
instance and the extensionality theorem \path{Perm.ext} allow consumers to
apply permutations, rewrite, and use \texttt{ext} without unfolding the
subgroup representation. The coercion to \texttt{Equiv.Perm A} is the
inherited subgroup coercion.

The finite-set property is \path{Perm.moved_finite}; the four equations in
Proposition~\ref{prop:subgroup} correspond respectively to
\path{Perm.moved_one}, \path{Perm.moved_mul_subset},
\path{Perm.moved_inv}, and \path{Perm.moved_conj}.

\subsection{Transpositions}
\label{sec:swaps}

Write $\swap a b$ for the permutation interchanging $a$ and $b$ and fixing
every other atom. An executable definition uses decidable equality on $A$;
classical reasoning can provide this locally when execution is not required.
The Lean operation is the qualified \texttt{Perm.swap}; it requires
\texttt{DecidableEq A}, but neither the permutation group nor its moved-point
laws require this assumption.

\begin{proposition}[Swap equations]
\label{prop:swaps}
For all $a,b,c\in A$,
\begin{gather*}
  \swap a b(a)=b,\qquad \swap a b(b)=a,\\
  c\ne a\ \land\ c\ne b\quad\Longrightarrow\quad\swap a b(c)=c,\\
  \swap a a=1,\qquad
  \swap a b^{-1}=\swap a b,\qquad
  \swap a b\swap a b=1,\\
  \pi\swap a b\pi^{-1}=\swap{\pi(a)}{\pi(b)}.
\end{gather*}
If $a\ne b$, then $\moved(\swap a b)=\{a,b\}$.
\end{proposition}

\begin{proof}
The defining cases show that applying a swap twice fixes each atom. They
also show that its moved-point set is contained in $\{a,b\}$, so it is a
finite-moved permutation. If $a\ne b$, both endpoints are moved, giving
the exact set. Conjugation sends $\pi(a)$ to $\pi(b)$ and $\pi(b)$ to
$\pi(a)$. A point outside those endpoints has inverse image outside
$\{a,b\}$ and is therefore fixed. Extensionality proves the conjugation
equation, including the case $a=b$.
\end{proof}

Within the namespace \texttt{Perm}, the endpoint equations are
\path{swap_apply_left} and \path{swap_apply_right}; the third-point equation
is \path{swap_apply_of_ne_of_ne}. The remaining equations are
\path{swap_self}, \path{swap_inv}, \path{swap_mul_self},
\path{conj_swap}, and \path{moved_swap}. The last theorem retains
the hypothesis $a\ne b$; a reflexive swap instead has empty moved-point set.
The implementation reuses Mathlib's corresponding equations for
\texttt{Equiv.swap}, including its conjugation equation, with the
finite-moved certificate proved when constructing \texttt{Perm.swap}.

\subsection{Controlled swap factorization}
\label{sec:controlled-swaps}

For an endpoint-pair list $l=[(a_1,b_1),\ldots,(a_n,b_n)]$, write
\[
  W(l)=\swap{a_1}{b_1}
    \bigl(\cdots(\swap{a_n}{b_n}1)\cdots\bigr),
  \qquad W([])=1.
\]
This denotes multiplication of permutations, so the rightmost swap acts
first. Prepending $(a,b)$ gives $W((a,b)::l)=\swap a b W(l)$.
The Lean expression is the ordinary list product
\begin{verbatim}
  (l.map (fun p => Perm.swap p.1 p.2)).prod
\end{verbatim}
No separate word datatype or evaluator is needed.

\begin{proposition}[Factorization within the original moved set]
\label{prop:controlled-factorization}
For every $\pi\in\Perm(A)$ there is a finite list $l$ of endpoint pairs
such that $W(l)=\pi$ and every $(a,b)\in l$ satisfies
\[
  a\ne b,\qquad a\in\moved(\pi),\qquad b\in\moved(\pi).
\]
\end{proposition}

\begin{proof}
Fix $M=\moved(\pi)$. In the ambient group of all bijections of $A$, let
\[
  G_M=\{\swap a b\mid a\ne b,\ a\in M,\ b\in M\}.
\]
Mathlib's restricted-transposition closure theorem characterizes membership
of $\sigma$ in the subgroup generated by $G_M$: its moved-point set must
be finite, and $\sigma(a)$ must lie in the orbit of $a$ under that subgroup
for every $a$. These conditions hold for $\pi$. Finiteness is part of
its certificate. If $\pi(a)=a$, the identity witnesses the orbit condition.
Otherwise $a\in M$ and $\pi(a)\in M$: an equality
$\pi(\pi(a))=\pi(a)$ would, by injectivity, imply $\pi(a)=a$.
Consequently $\swap a{\pi(a)}$ is a generator sending $a$ to $\pi(a)$.
Thus the underlying bijection of $\pi$ lies in the closure of $G_M$.

Left closure induction extracts the desired list. The identity uses the
empty list. Left multiplication by a generator prepends its endpoint pair;
left multiplication by the inverse of a generator prepends the same pair,
because a swap is its own inverse. Throughout this induction, every endpoint
belongs to the fixed set $M$, the moved set of the original $\pi$.
The induction records that the underlying bijection of $W(l)$ equals the
current ambient permutation. The existing identity and multiplication
equations for \texttt{toEquiv} preserve this equality at each step.
Subtype extensionality then yields $W(l)=\pi$ in $\Perm(A)$.
\end{proof}

The finite-set reduction behind the adopted Mathlib closure theorem explains
why this argument terminates. For $\sigma\in\Perm(A)$ and a moved atom $a$, put
$\tau=\swap a{\sigma(a)}$ and $\rho=\tau\sigma$. Then $\rho(a)=a$ and
\[
  \moved(\rho)\subseteq\moved(\sigma)\setminus\{a\},
  \qquad |\moved(\rho)|<|\moved(\sigma)|,
  \qquad \sigma=\tau\rho.
\]
The inclusion says that a point moved after this correction was already
moved and differs from $a$. Since $a$ was moved, the inclusion gives a
strict decrease in finite cardinality. It does not assert exact deletion
of just one point: if $\sigma$ is a two-cycle, $\rho=1$ and both endpoints
disappear. The package reuses Mathlib's finite-set induction rather than
copying that reduction. The local restriction to $G_M$ and the subsequent
closure induction supply the original-set endpoint bound.

The declaration is \path{Perm.swap_factorization}. Its source is
\begin{quote}
\small\path{Package/Foundations/SwapFactorization.lean}.
\end{quote}
The proof uses \path{mem_closure_isSwap} and
\path{Subgroup.closure_induction_left}. Their sources are respectively
\begin{quote}
\small
\path{Mathlib/GroupTheory/Perm/ClosureSwap.lean},\\
\path{Mathlib/Algebra/Group/Subgroup/Pointwise.lean},
\end{quote}
at the pinned revision recorded in Section~\ref{sec:evidence}.
The empty-product certificate is definitional, and
\path{Perm.toEquiv_mul} tracks each prepended factor; the proof introduces
no separate product homomorphism or public helper.
In particular, the proof uses a restricted generating set, rather than
assuming that unrestricted swap generation already provides endpoint control.

\begin{proposition}[Factorization avoiding a pointwise fixed set]
\label{prop:avoidance-factorization}
Let $S\subseteq A$ be any set, and suppose $\pi(a)=a$ for every $a\in S$.
There is a finite list $l$ with $W(l)=\pi$ such that every $(a,b)\in l$
satisfies
\[
  a\ne b,\qquad a\notin S,\qquad b\notin S.
\]
\end{proposition}

\begin{proof}
Every member of $\moved(\pi)$ lies outside $S$, since a member of $S$
is fixed. Apply Proposition~\ref{prop:controlled-factorization} and use
this fact at each endpoint.
\end{proof}

The declaration is \path{Perm.swap_factorization_avoiding}.
Both propositions retain \texttt{DecidableEq A} for the swap operation,
which classical reasoning can supply locally. They require no infinitude,
countability, finiteness, or inhabitant assumption on $A$. The avoidance
parameter is an arbitrary \texttt{Set A}; neither $S$ nor its complement
has to be finite. If $\pi=1$, endpoint membership in its empty moved set
forces the controlled list to be empty. This also covers empty and singleton
atom carriers, where every permutation is the identity. The results give
existential factor lists, with no canonical choice, minimal length, uniqueness,
or executable factoring algorithm promised.

\subsection{Selected actions and canonical data shapes}
\label{sec:actions}

An action on $X$ consists of an operation $\pi\act x$ satisfying
\[
  1\act x=x,\qquad (\pi\sigma)\act x=\pi\act(\sigma\act x).
\]
The interface uses \texttt{MulAction (Perm A) X} directly. No additional
action class or independently competing scalar operation is installed.
Standard group-action laws give
\[
  \pi^{-1}\act(\pi\act x)=x,\qquad
  \pi\act(\pi^{-1}\act x)=x,\qquad
  \pi\act x=\pi\act y\ \Longleftrightarrow\ x=y.
\]
These equations concern the selected action; a different action on the same
carrier is different input to subsequent theorems.

The canonical action on atoms is evaluation, $\pi\act a=\pi(a)$,
as stated by \path{Perm.smul_atom}. It is obtained by restricting Mathlib's
\path{Equiv.Perm.applyMulAction} to the subgroup \texttt{permSubgroup A}.
Thus the computation theorem describes an inherited action; it does not
introduce another action instance on atoms.
The standard product action is componentwise:
\[
  \pi\act(x,y)=(\pi\act x,\pi\act y).
\]
Both projections therefore commute with renaming. This construction permits
the two factors, as well as the atom carrier, to inhabit different universes.

For $S\in\Finset(A)$, the selected action is direct image,
$\pi\act S=\pi[S]$. Its public equations include
\begin{align*}
  a\in\pi\act S&\quad\Longleftrightarrow\quad\pi^{-1}(a)\in S,\\
  \pi\act\varnothing&=\varnothing,\\
  \pi\act\{a\}&=\{\pi(a)\},\\
  \pi\act(S\cup T)&=(\pi\act S)\cup(\pi\act T).
\end{align*}
Membership follows from the inverse bijection; the other identities follow
from the corresponding direct-image identities. In Lean this action reuses
Mathlib's scoped pointwise action on finite sets. A consumer must explicitly
write \texttt{open scoped Pointwise}; opening the scope in a dependency is
not a promise that it is active downstream. Finite-set computations may use
local classical decidable equality. No overlapping global finite-set action
is needed. The image and membership equations are \path{Perm.smul_finset}
and \path{Perm.mem_smul_finset}. The standard action laws and the product,
empty-set, singleton, and union equations are supplied by Mathlib; the package
does not duplicate them under fresh names.

\subsection{Swap invariance for a selected action}
\label{sec:swap-invariance}

Fix an arbitrary selected action of $\Perm(A)$ on $X$. For any set
$S\subseteq A$ and value $x\in X$, invariance under the pointwise fixer
of $S$ can be tested using only swaps outside $S$.

\begin{proposition}[Pointwise-fixer invariance and outside swaps]
\label{prop:swap-invariance}
For an arbitrary set $S\subseteq A$ and an arbitrary $x\in X$,
\begin{align*}
  &\bigl(\forall\pi\in\Perm(A),\
      (\forall c\in S,\ \pi(c)=c)\Longrightarrow\pi\act x=x\bigr)\\
  &\hspace{2em}\Longleftrightarrow
    \bigl(\forall a,b\in A,\
      a\notin S\Longrightarrow b\notin S
      \Longrightarrow\swap a b\act x=x\bigr).
\end{align*}
\end{proposition}

\begin{proof}
Suppose every swap with endpoints outside $S$ fixes $x$, and let $\pi$
fix $S$ pointwise. Proposition~\ref{prop:avoidance-factorization} gives
a list $l$ of such swaps whose product is $\pi$. Induct on that list.
The empty product fixes $x$ by the identity-action law. For a prepended
pair $(a,b)$, the product-action law gives
\[
  W((a,b)::l)\act x
    =\swap a b\act(W(l)\act x)
    =\swap a b\act x=x.
\]
Thus $\pi\act x=x$. Conversely, if $a,b\notin S$, every $c\in S$
differs from both endpoints, so $\swap a b(c)=c$ by the third-point swap
equation. Apply the pointwise-fixer hypothesis to this swap.
\end{proof}

The declarations are \path{Perm.smul_eq_of_swap_smul_eq} for
the implication from outside-swap invariance to pointwise-fixer invariance,
and \path{Perm.forall_smul_eq_iff_swap_smul_eq} for the equivalence.
They use only \texttt{DecidableEq A} and
\texttt{MulAction (Perm A) X}, with independent atom and carrier universes.
The swap premise permits equal endpoints; those cases follow from
$\swap a a=1$ and $1\act x=x$, so requiring only distinct-endpoint swaps
would be equivalent.

Pointwise fixation here means $\pi(c)=c$ for each $c\in S$; setwise
preservation $\pi[S]=S$ is a weaker condition. The finite moved-point set
of a permutation also differs from support of a value in an arbitrary
action. No assumption that $x$, or every value of $X$, has finite support
enters this proposition. It provides the permutation input for the later
support layer, \texttt{PKG-F04}; a support predicate, least support, and
freshness remain future work.

\subsection{Explicit discrete data}
\label{sec:discrete}

\begin{definition}[Discrete permutation set]
For a group $G$ and a set $X$, the \emph{discrete $G$-set on $X$} has
underlying set $X$ and action $g\act x=x$ for every $g\in G$ and $x\in X$.
For $G=\Perm(A)$ we denote this construction by $\Delta_A(X)$.
\end{definition}

This is the mathematical construction encoded by \texttt{Discrete A X}.
It is the discrete group action of Pitts~\cite[Example~1.5, p.~15]{pitts2013}
and, for an infinite atom set, the discrete nominal set of
\cite[Example~2.5, p.~30]{pitts2013}. The implementation uses a canonically
bijective copy $\Discrete_A(X)$ of the carrier, with constructor
$\operatorname{mk}$ and projection $\operatorname{val}$, and action
\[
  \pi\act\operatorname{mk}(x)=\operatorname{mk}(x).
\]

Thus $\operatorname{val}(\operatorname{mk}(x))=x$,
$\operatorname{mk}(\operatorname{val}(d))=d$, and
$\operatorname{val}(\pi\act d)=\operatorname{val}(d)$.
The structure \texttt{Discrete A X} makes the trivial-action
choice explicit, even when $X=A$. It supports arbitrary type universes,
including empty and singleton data represented by universe-correct
\texttt{PEmpty} and \texttt{PUnit}. More precisely, if
$A:\texttt{Type}\ u$ and $X:\texttt{Type}\ v$, then
$\Discrete_A(X):\texttt{Type}\ v$: the atom type is a phantom parameter,
not stored data. This wrapper carries the only new action instance in the
action module.

In words, $\Discrete_A(X)$ means ``values of $X$ that are unaffected by
renaming atoms of $A$.'' The parameter $A$ specifies the acting group;
$X$ specifies the data being carried. For example, $\Discrete_A(\mathbb N)$
can represent numerical annotations that stay unchanged during renaming.
The construction also permits $X=A$: the wrapped datum is then treated as
a fixed label, as detailed in Example~\ref{ex:discrete}. This choice does
not assert that an action already supplied on $X$ is trivial.
If $X$ already has an action, the construction forgets that action and
equips its underlying data with the trivial one. It does not take the
subset of fixed points or the quotient into orbits.

Every orbit is a singleton. In the usual support terminology, the empty
set supports each wrapped value: every permutation fixes it, including
every permutation fixing the empty set pointwise. Since the empty set is
contained in every other set, it is its least support. This is a
mathematical consequence of the trivial action; the package's general
support and freshness interfaces remain later work.

\paragraph{Why a single-field structure?}
The complete data declaration is
\begin{verbatim}
structure Discrete (A : Type u) (X : Type v) : Type v where
  val : X
\end{verbatim}
Only the underlying value must be stored. The name \texttt{val} is a
conventional projection name with no mathematical significance; it could
equally be called \texttt{value}. A single field does not mean a single
inhabitant: each element of $X$ determines one wrapped value. The trivial
action and its laws are supplied uniformly by a typeclass instance, so
individual values need not store an action, an atom, or a support certificate.

The structure makes $\Discrete_A(X)$ a separate Lean type, not
definitionally equal to $X$. Its constructor and projection explicitly
cross this boundary. This lets the canonical atom action on $A$ and the
trivial action on $\Discrete_A(A)$ coexist under ordinary instance search.

\paragraph{Alternative declarations.}
An abbreviation \texttt{abbrev Discrete A X := X} unfolds during reducible
unification. It therefore does not supply the distinct carrier boundary
chosen here for selecting a different action. A regular definition
\texttt{def Discrete A X := X} is a viable type-tag alternative: it is
still definitionally equal to $X$, while its unfolding behavior differs
from that of an abbreviation. Mathlib uses this technique for
\texttt{Additive} and \texttt{Multiplicative} in
\path{Mathlib/Algebra/Group/TypeTags/Basic.lean}, inspected at the pinned
Mathlib revision recorded in Section~\ref{sec:evidence}.
The pinned Lean 4.34.1 instance search does not unfold an ordinary
semireducible definition. Thus a regular definition can separate the
tagged and underlying action instances; a structure is not required for
that purpose alone. This agrees with the general account in the
\href{https://lean-lang.org/doc/reference/latest/Type-Classes/Instance-Synthesis/}{Lean reference on instance synthesis};
the installed compiler sources and the probes below determine the
version-specific claims here.

The definition approach permits direct reuse of underlying values,
functions, predicates and containers through definitional equality. For
example, a list of $X$ can also be checked as a list of tagged $X$ without
mapping a constructor over it. With the structure, that passage is explicit.
Both approaches already have reflexive conversion round trips. Neither
automatically supplies every underlying typeclass instance: our probes
needed explicit forwarding of \texttt{Inhabited} and \texttt{DecidableEq}
for a regular tag, just as the current structure needs corresponding
instances when a client requires them.

\paragraph{What the elaboration probes distinguish.}
Action selection depends on the type Lean infers for an expression, not
only on a type annotation that it checks against. In a scratch regular-def
tag \texttt{D}, with the trivial action registered, a variable
\texttt{d : D Bool Bool} is fixed. Yet an existing raw Boolean value checked
against that definitionally equal type can still select the raw atom action.
For the swap $\tau=\swap{\mathrm{false}}{\mathrm{true}}$, the tested expression
\begin{verbatim}
  tau • (false : D Bool Bool)
\end{verbatim}
elaborates to the raw action and equals \texttt{true}. A named identity-valued
constructor \texttt{D.mk false}, whose declared result is the tagged type,
instead selects the trivial action and stays fixed. Named equivalence-based
conversions and explicitly typed local binders also worked. The Boolean
carrier deliberately tests the action layer without any infinitude premise.

The same distinction appeared for an existing raw pair and finite set:
checking either against a definitionally equal tagged container type did
not necessarily change its inferred action. A named conversion with the
tagged result type did. Reusing an already elaborated function likewise
preserves the operations and instances selected in its body. These are
elaboration choices, not inconsistent proofs or a mathematical failure of
the discrete construction. They explain why a regular-def interface still
needs an explicit conversion discipline for predictable action selection.

Another option is to define the trivial action directly on $X$ and select
it locally or pass it explicitly. That avoids a wrapper, but simultaneous
reasoning with two actions on the same carrier then requires managing the
two structures explicitly. The present single-field structure chooses a
stronger type boundary at the cost of constructor/projection conversions.
This is an interface decision, not a mathematical necessity or a claim
that a regular-definition implementation cannot work.

\paragraph{Mathlib reuse and the resulting choice.}
Mathlib provides precedents for both encodings. Its \texttt{OrderDual}
tag explicitly recommends conversion functions even though the
opposite-order carrier is definitionally equal to the original.
The source is \path{Mathlib/Order/OrderDual.lean}.
Conversely, the categorical \texttt{Discrete} chooses a one-field structure
to control passage to the underlying type.\footnote{See
\path{Mathlib/CategoryTheory/Discrete/Basic.lean}.}
That declaration encodes a \emph{discrete category}, not a discrete
permutation set, so its name is not a reason to reuse it as our nominal carrier.

The implementation states the discrete action directly:
\begin{verbatim}
instance : MulAction (Perm A) (Discrete A X) where
  smul _ d := d
  one_smul _ := rfl
  mul_smul _ _ _ := rfl
\end{verbatim}
The first line of the record ignores the permutation and returns the same
value. The identity permutation therefore fixes every value. Applying two
permutations successively and applying their product both return that same
value, which explains the two \texttt{rfl} proofs. The record supplies the
standard Mathlib action interface with the operation we want.

We checked Mathlib's alternative action constructors and verified that one
produces a definitionally equal record. The explicit form above is retained
for readability: it displays the mathematical action and both laws immediately.
The alternatives and their exact source references remain in the research
record. This choice adds no new general theory of group actions.

We retain the present structure: for this action-facing API,
explicit carrier boundaries currently outweigh the definition's container
conversion advantage. This conclusion can be revisited when a real client
demonstrates that conversion costs dominate; it does not prescribe
structures for all future type tags.

The ordinary type equivalence between $\Discrete_A(X)$ and $X$ does not
identify their actions. If $X$ already carries an independently selected
action, its projection is equivariant exactly when every $x\in X$ is fixed
by every permutation. No automatic discrete action on arbitrary external
carriers is part of the interface.
The operation is \path{Discrete.equiv}, with computation equations
\path{Discrete.equiv_apply} and \path{Discrete.equiv_symm_apply}.
The trivial-action equations are \path{Discrete.smul_mk} and
\path{Discrete.smul_val}.

\subsection{Equivariance of ordinary functions}
\label{sec:equivariance}

\begin{definition}
Given selected actions on $X$ and $Y$, an ordinary function $f:X\to Y$
is \emph{equivariant} when
\[
  \Equivariant_A(f)\quad\Longleftrightarrow\quad
  \forall\pi\in\Perm(A)\ \forall x\in X,\quad
  f(\pi\act x)=\pi\act f(x).
\]
\end{definition}

In Lean, this unbundled property is written \texttt{Equivariant A f},
with the atom type $A$ an explicit input.
It is not an \texttt{outParam} inferred from the carrier. The domain and
codomain have independently selected \texttt{MulAction} instances and may
inhabit different universes. Defining equivariance requires neither a
supported-function wrapper nor an action on the function type itself.
No global action on truth values is introduced.

\begin{proposition}[Elementary equivariant maps]
\label{prop:equivariant}
Identity maps and product projections are equivariant. If $f:X\to Y$
and $g:Y\to Z$ are equivariant, so is $g\circ f$. If $f:X\to Y$ and
$h:X\to Z$ are equivariant, so is $x\mapsto(f(x),h(x))$.
\end{proposition}

\begin{proof}
Identity is immediate, and the projection statements follow from the
componentwise product action. For composition,
\[
  (g\circ f)(\pi\act x)
  =g(\pi\act f(x))
  =\pi\act g(f(x)).
\]
For pairing, apply the two equivariance equations componentwise.
\end{proof}

The corresponding declarations are \path{Equivariant.id},
\path{Equivariant.fst}, \path{Equivariant.snd},
\path{Equivariant.comp}, and \path{Equivariant.pair}.
Composition takes its certificates in the order
\texttt{Equivariant.comp hg hf}, yielding equivariance of $g\circ f$.
Pairing takes the certificates for the first and second components in that
order. The proofs use ordinary function application and the selected actions;
no function bundle is needed.

Later supported-map certificates can weaken equivariance by restricting
the permutations to those fixing a finite set of parameters. Fixing a
parameter in a jointly equivariant operation need not leave an equivariant
function. No finite-support closure theorem is claimed in this increment.

\subsection{Three distinctions between actions}
\label{sec:counterexamples}

The following examples require just two distinct atoms $a\ne b$. Put
$\tau=\swap a b$, so $\tau(a)=b$ and $\tau\ne1$. They isolate choices
of action rather than difficulties concerning infinitude or support.

\begin{example}[An atom and a discrete copy]
\label{ex:discrete}
The canonical action moves $a$, since $\tau\act a=b\ne a$.
The discrete action fixes $\operatorname{mk}(a)$, since
$\tau\act\operatorname{mk}(a)=\operatorname{mk}(a)$.
Consequently, $\operatorname{val}:\Discrete_A(A)\to A$ is a bijection of
types but is not equivariant for these actions: applying it after renaming
gives $a$, whereas renaming its value gives $b$.
\end{example}

\begin{example}[Left multiplication and conjugation]
\label{ex:group-actions}
A group acts on itself by left multiplication, $\pi\act\sigma=\pi\sigma$.
It also admits the conjugation action $\pi\star\sigma=\pi\sigma\pi^{-1}$.
These disagree at the identity:
\[
  \tau\act1=\tau\ne1,\qquad \tau\star1=1.
\]
The bare group $\Perm(A)$ retains its standard left action. The moved-point
conjugation law in Proposition~\ref{prop:subgroup} is a group equation;
it does not install conjugation as that group's scalar action.
\end{example}

\begin{example}[Pointwise and conjugation actions on functions]
\label{ex:function-actions}
For ordinary $f:X\to Y$, the pointwise action uses only the codomain action:
\[
  (\pi\act_{\mathrm{pt}}f)(x)=\pi\act f(x).
\]
When the domain also has an action, the conjugation formula is
\[
  (\pi\act_{\mathrm{conj}}f)(x)
    =\pi\act f(\pi^{-1}\act x).
\]
On the identity function of the canonical atom carrier these differ:
\[
  (\tau\act_{\mathrm{pt}}\id_A)(a)=b,\qquad
  (\tau\act_{\mathrm{conj}}\id_A)(a)=a.
\]
Equivariance is precisely being fixed by every conjugation operation. Indeed,
substituting $\pi^{-1}\act x$ in its defining equation gives
$\pi\act f(\pi^{-1}\act x)=f(x)$; conversely, substitute $\pi\act x$
in this fixed-point equation. It is not the fixed-point condition for the
ordinary pointwise action. The interface does not replace Mathlib's
pointwise action on bare function types or infer that arbitrary functions
are finitely supported.
\end{example}

These are manuscript calculations explaining the selected actions. Temporary
Lean checks confirmed the three distinctions through the public import and
checked that a structure without a declared action does not acquire one
automatically. They do not introduce additional public theorems or a
persistent example library. The absence of an inferred action is an interface
property, not a claim that no mathematical action can be constructed.

\subsection{Implementation correspondence and evidence}
\label{sec:evidence}

The foundation contract is divided as follows. The permutation,
factorization, and action modules have been checked, together with their
source coverage and axiom dependencies.
The declaration references identify the implemented statements; the manuscript
also explains consequences that are not exported as separate named theorems.

\begin{center}
\small
\begin{tabular}{@{}>{\raggedright\arraybackslash}p{0.37\linewidth}
  >{\raggedright\arraybackslash}p{0.57\linewidth}@{}}
\hline
Contract & File \\
\hline
Subgroup, application, moved sets, swaps &
  \path{Package/Foundations/Permutation.lean} \\
Controlled and avoiding factorization &
  \path{Package/Foundations/SwapFactorization.lean} \\
Actions, swap criterion, discrete data, equivariance &
  \path{Package/Foundations/Action.lean} \\
Public import & \path{Package.lean} \\
Module-origin axiom audit & \path{Package/Tests/AxiomAudit.lean} \\
Import and source-coverage check & \path{Package/Scripts/check-imports.py} \\
\hline
\end{tabular}
\end{center}

The implementation uses the repository's pinned Lean and Mathlib 4.34.1
versions. The Mathlib revision, agreeing with the manifest and installed
checkout, is
\begin{quote}
\small\texttt{d13f23b723b8a846827a245b89c10fc7d3f11612}.
\end{quote}
The \texttt{F03b} checks concern an uncommitted increment. Its branch is
\path{fasapa/nominal-package}, and its committed \texttt{F02/F03a} base is
\begin{quote}
\small\texttt{3d2196a57b8964084dd58a65a0cb1c0be50b1592}.
\end{quote}
The final verification commands were
\begin{quote}
\small
\begin{verbatim}
lake build Package +Package.Tests.AxiomAudit
lake env lean \
  /tmp/nominal-f03b-execution/FactorizationContracts.lean
lake env lean \
  /tmp/nominal-f03b-execution/ActionContracts.lean
lake env lean \
  /tmp/nominal-f02-f03a-execution/ActionContracts.lean
lake env lean \
  /tmp/nominal-f02-f03a-execution/NegativeContracts.lean
lake env lean Package/Tests/AxiomAudit.lean
python3 Package/Scripts/check-imports.py
git diff --check
\end{verbatim}
\end{quote}
All commands exited successfully. The direct audit and temporary consumers
produced no warnings. Iterative builds rebuilt the changed package modules
using cached dependencies. The final combined build completed 967 jobs,
rebuilding the audit and reusing those checked modules.
This evidence does not claim a fresh project build, a dependency bootstrap,
or a rebuild of the unchanged reference development.

The source-coverage checker reports four production source modules: the public
root and the three foundation modules. The audit category contains one module,
whose import closure reaches all five Lean source files. The checker validates
headers using the pinned Lean parser, rejects missing or unclassified package
sources and forbidden dependencies, and checks reachability from the respective
roots. Its thirteen self-tests passed during the preceding \texttt{F02/F03a}
increment; the checker is unchanged and those self-tests were not rerun for
\texttt{F03b}. In particular, the production closure
has no dependency on the earlier nominal library, its clients, or research
probes. Header discovery complements the full Lean build; it does not replace
proof elaboration or checking of module bodies.

The axiom audit examines declarations by their originating modules, including
private and generated declarations, rather than selecting a namespace prefix.
It checked 81 production declarations from the three foundation modules;
the public umbrella introduces no declarations of its own. Transitive axiom
dependencies were confined to \texttt{propext}, \texttt{Classical.choice}, and
\texttt{Quot.sound}. The audit rejects a count of zero or any additional axiom,
and checks its rejection logic with simulated axiom names rather than new
axiom declarations. Its direct invocation repeats the traversal even when
Lake can reuse a compiled audit module.

The audit prints dependencies for all four new results:
\begin{quote}
\small
\path{Perm.swap_factorization},\\
\path{Perm.swap_factorization_avoiding},\\
\path{Perm.smul_eq_of_swap_smul_eq},\\
\path{Perm.forall_smul_eq_iff_swap_smul_eq}.
\end{quote}
It also retains representative checks for
\path{Perm.moved_conj}, \path{Perm.conj_swap}, \path{Perm.smul_atom},
\path{Discrete.smul_val}, and \path{Equivariant.comp}; these likewise
contain only the permitted standard axioms. These results establish the
permutation, factorization, and action layer described here. They do not supply finite
support, least support, predicates, or the future case studies.

The new temporary consumers use only \texttt{import Package}. They check
that all factors fix an atom fixed by the original permutation; identity,
empty, and singleton carriers give empty controlled lists; and a nontrivial
swap on a finite carrier gives a nonempty list. A noncommuting three-point
product distinguishes the rightmost-first order from its reverse. Avoidance
is also exercised on the infinite set $\mathbb N\setminus\{0,1\}$, without
restricting the theorem's arbitrary-set parameter.

The action consumers recover both component equalities from invariance of
a product with independent carrier universes, exercise both directions of the criterion at
$S=\varnothing$, and handle $S=A$ and a premise restricted to distinct
endpoints. Rerun \texttt{F03a} consumers preserve the atom, product,
finite-set, discrete, bare-group, and pointwise-function action distinctions
and the expected missing-action diagnostic. These successful checks are
scratch evidence rather than maintained package modules; persistent usage
examples belong to the future case studies.

The manuscript was compiled from \path{docs/article/} with
\begin{quote}
\small
\begin{verbatim}
latexmk -pdf -interaction=nonstopmode -halt-on-error \
  -outdir=/tmp/nominal-package-article-build main.tex
\end{verbatim}
\end{quote}
Generated files remain outside the source tree. The final log was checked
for unresolved references and layout diagnostics, and the changed pages'
extracted text was inspected. PDF compilation checks the document artifact;
the separate source correspondence check verifies the four Lean signatures,
their assumptions, product order, and proof route.

\section{Finite support calculus}
\label{sec:finite-support}

Fix an atom carrier $A$ and selected actions of $\Perm(A)$ on $X$ and $Y$.
The atom action is canonical, $\pi\act a=\pi(a)$, and the universes of
all three carriers are independent. Finite support is a property of an
individual element; it does not require every element of its carrier to
have finite support. Infinitude enters the calculus only at intersection.

\subsection{Supporting bounds}
\label{sec:finite-support-definition}

\begin{definition}
A finite set $S\subseteq A$ \emph{supports} $x\in X$ when
\[
  \operatorname{Supports}(S,x)
  \quad\Longleftrightarrow\quad
  \forall\pi\in\Perm(A),\quad
  \bigl(\forall a\in S,\ \pi(a)=a\bigr)
  \Longrightarrow \pi\act x=x.
\]
The element $x$ is \emph{finitely supported} if such a finite $S$ exists.
\end{definition}

The definition makes a supporting set a sufficient bound, without choosing
one or asserting minimality. It requires pointwise fixation of $S$;
setwise preservation may permute its members. Nor does fixation of $x$
imply pointwise fixation of a supporting set. That reverse implication
belongs to the stronger notion of strong support. Support also depends on
the action: for example, support on bare $\Perm(A)$ refers to left
multiplication, not silently to conjugation
(Example~\ref{ex:group-actions}).

The Lean predicate \texttt{Supports S x}, with \texttt{S : Finset A},
is an abbreviation for Mathlib's two-carrier predicate
\begin{verbatim}
  MulAction.Supports (Perm A) (S : Set A) x
\end{verbatim}
This reuses its pointwise-fixer semantics with the canonical action on
atoms and the selected action on values. The atom sort is determined by
$S$, whereas \texttt{FinitelySupported A x} takes it explicitly.
Both predicates inhabit \texttt{Prop}, so neither supplies an observable
choice of bound. Neither definition requires decidable equality
or infinitude of $A$.

\subsection{Elementary bounds and equivariant images}
\label{sec:finite-support-basic}
\label{sec:finite-support-products}

\begin{proposition}[Elementary support bounds]
\label{prop:support-basic}
\label{prop:support-products}
\label{prop:finitely-supported-basic}
For finite bounds $S,T\subseteq A$:
\begin{enumerate}
\item If $S$ supports $x$ and $S\subseteq T$, then $T$ supports $x$.
\item The empty set supports $x$ exactly when every permutation fixes $x$.
\item If $f:X\to Y$ is equivariant and $S$ supports $x$, then $S$ supports
  $f(x)$.
\item Under the componentwise action, $S$ supports $(x,y)$ exactly when it
  supports both components. If $S$ supports $x$ and $T$ supports $y$,
  then $S\cup T$ supports $(x,y)$.
\end{enumerate}
Consequently, equivariant images and products preserve finite supportedness.
None of these conclusions requires infinitude.
\end{proposition}

\begin{proof}
A permutation fixing $T$ pointwise also fixes $S\subseteq T$; for the
empty set the pointwise condition is vacuous. Equivariance and support give
$\pi\act f(x)=f(\pi\act x)=f(x)$.
For products, fixation is componentwise; enlarge separate bounds to
$S\cup T$. Existential witnesses for the image and product are the same
bound and the union, respectively.
\end{proof}

The image theorem is a bound, not an equality of least supports. For
example, a constant map into $\Discrete_A(\mathbb B)$ is equivariant.
Its value inherits the bound $\{a\}$ of an atom $a$, but also has the
strictly smaller empty bound. The image theorem needs only the
ordinary-function predicate $\Equivariant_A(f)$, without an action or a
support certificate on the function itself.

\subsection{Support transport by conjugation}
\label{sec:support-transport}

\begin{proposition}[Conjugation transport]
\label{prop:support-transport}
For every $\pi\in\Perm(A)$ and finite $S\subseteq A$,
\[
  \operatorname{Supports}(\pi[S],\pi\act x)
  \quad\Longleftrightarrow\quad
  \operatorname{Supports}(S,x).
\]
No infinitude or commuting-action hypothesis is required.
\end{proposition}

\begin{proof}
Suppose $S$ supports $x$ and $\sigma$ fixes $\pi[S]$ pointwise.
For $a\in S$, one has $\sigma(\pi(a))=\pi(a)$, so
$(\pi^{-1}\sigma\pi)(a)=a$.
Thus $\pi^{-1}\sigma\pi$ fixes $S$ pointwise and hence fixes $x$.
The action laws give
\[
  \sigma\act(\pi\act x)
  =\pi\act\bigl((\pi^{-1}\sigma\pi)\act x\bigr)
  =\pi\act x.
\]
For the converse, apply this implication to $\pi^{-1}$, the bound
$\pi[S]$ and the element $\pi\act x$, then cancel the inverse actions.
\end{proof}

The formal statement is \path{supports_smul_iff}. Conjugation is essential
to this proof: renaming transports the pointwise fixer of $S$ to that of
$\pi[S]$. Mathlib's \path{MulAction.Supports.smul} instead assumes
commuting scalar actions on both carriers, a condition that general
permutations do not satisfy. The conjugation argument accommodates their
noncommutativity directly.

\begin{proposition}[Existential transport]
\label{prop:finitely-supported-transport}
For every $\pi\in\Perm(A)$,
\[
  \operatorname{FinitelySupported}_A(\pi\act x)
  \quad\Longleftrightarrow\quad
  \operatorname{FinitelySupported}_A(x).
\]
\end{proposition}

Transport sends a witness $S$ to $\pi[S]$; applying $\pi^{-1}$ gives the
converse. Thus finite supportedness is
invariant under renaming even when no bound is selected as data.

\subsection{Swaps and intersection}
\label{sec:support-swap-criterion}
\label{sec:support-intersection}

\begin{proposition}[Support tested by swaps]
\label{prop:support-swap-criterion}
A finite $S\subseteq A$ supports $x$ if and only if
\[
  \forall a,b\in A,\quad a\notin S\ \land\ b\notin S
  \Longrightarrow\swap a b\act x=x.
\]
Equal endpoints are allowed, and no infinitude premise is required.
\end{proposition}

This specializes Proposition~\ref{prop:swap-invariance} to a finite bound.
It turns a quantification
over all finite permutations into a local condition on two atoms.
Intersection then reduces to finding a spare atom joining two such local
conditions.

\begin{proposition}[Intersection of finite support bounds]
\label{prop:finite-support-intersection}
Assume $A$ is infinite. If finite $S,T\subseteq A$ both support $x$,
then $S\cap T$ supports $x$.
\end{proposition}

\begin{proof}
By the swap criterion, it suffices to show that $\swap a b$ fixes $x$
whenever $a,b\notin S\cap T$. The case $a=b$ is immediate. Otherwise,
choose
\[
  c\notin S\cup T\cup\{a,b\}.
\]
Such an atom exists because the excluded set is finite. The atom $a$ lies
outside at least one of $S,T$, and $c$ lies outside both; therefore
$\swap a c$ fixes $x$. Similarly, $\swap c b$ fixes $x$. Conjugation gives
\[
  \swap a c\swap c b\swap a c
  =\swap{\swap a c(c)}{\swap a c(b)}
  =\swap a b.
\]
All three factors on the left fix $x$, so their product does too.
\end{proof}

In \path{supports_inter}, infinitude is used only to obtain $c$ by
Mathlib's \path{Finset.exists_notMem}. There is no assumption that all
values of $X$ are finitely supported. The role of the atom hypothesis is
mathematical, as the following example shows.

\begin{example}[Intersection can fail on finite atoms]
\label{ex:finite-atom-intersection}
Let $A=\mathbb B$ carry its canonical action, and take
$x=\mathrm{false}$. Both $\{\mathrm{false}\}$ and
$\{\mathrm{true}\}$ support $x$, but their empty intersection does not.
\end{example}

\begin{proof}
The first bound is immediate. A permutation fixing $\mathrm{true}$ also
fixes $\mathrm{false}$: injectivity prevents it from sending both atoms to
$\mathrm{true}$, and these are the only atoms. This gives the second
bound. Yet $\swap{\mathrm{false}}{\mathrm{true}}$ moves $x$, so the empty
set cannot support it.
\end{proof}

The example uses the same finite-permutation group and support definition
as the infinite-atom theorem. It shows why a sufficient support bound
cannot be identified with a least bound and why intersection needs an
atom hypothesis.

\subsection{Assumptions and support witnesses}
\label{sec:support-assumptions}

The mathematical distinction between a supplied finite bound and existence
of a bound has a useful consequence for the Lean interface. Visible union,
image and intersection expressions require decidable equality. When such
a set is constructed only as an existential witness, classical equality
can remain local to the proof: finite supportedness still carries no
computational choice of bound. The assumptions therefore separate as follows.

\begin{center}
\small
\begin{tabular}{@{}>{\raggedright\arraybackslash}p{0.66\linewidth}
  >{\raggedright\arraybackslash}p{0.28\linewidth}@{}}
\hline
Statements & Atom assumptions \\
\hline
Definitions, enlargement, empty support, equivariant images,
common product bounds, closure and transport of finite supportedness & None \\
Visible finite-set image and union bounds, swap criterion &
  \texttt{DecidableEq A} \\
Intersection of finite support bounds &
  \texttt{DecidableEq A}, \texttt{Infinite A} \\
\hline
\end{tabular}
\end{center}

All rows use the selected actions, with independent carrier universes.
Countability, carrier-wide supportedness and commuting-action assumptions
are unnecessary. The calculus therefore applies to finitely supported
elements within arbitrary permutation actions, before imposing any
condition on the carrier as a whole.

\subsection{Nominality and least support}
\label{sec:least-support}

\begin{definition}[Nominality certificate]
A selected action on $X$ is \emph{nominal} when every $x\in X$ is finitely
supported. The certificate asserts
\[
  \forall x\in X,\quad
  \exists S\in\Finset(A),\quad\operatorname{Supports}(S,x).
\]
\end{definition}

This condition is meaningful for any atom carrier. In Lean,
\texttt{Nominal A X} is a class in \texttt{Prop}, parameterized by the
already selected \texttt{MulAction (Perm A) X}. Its sole field proves finite
supportedness of every element. It contains neither an action nor a chosen
finite set, and requires neither infinitude nor decidable equality. Atom and
carrier universes remain independent; the atom parameter is explicit.

Least support concerns an individual element. Write
$\operatorname{supp}_A(x)$ for its least finite supporting set when this
exists. The next result requires finite supportedness of $x$, without a
nominality certificate for all of $X$.

\begin{proposition}[Unique least finite support]
\label{prop:least-support}
Assume $A$ is infinite and $x\in X$ is finitely supported. There is a unique
finite set $L\subseteq A$ such that
\[
  \operatorname{Supports}(L,x)
  \quad\text{and}\quad
  \forall T\in\Finset(A),\quad
  \operatorname{Supports}(T,x)\Longrightarrow L\subseteq T.
\]
Consequently, for every finite $S\subseteq A$,
\[
  \operatorname{Supports}(S,x)
  \quad\Longleftrightarrow\quad\operatorname{supp}_A(x)\subseteq S.
\]
\end{proposition}

\begin{proof}
Strict inclusion of finite sets is well founded. Since $x$ has a finite
support, choose an inclusion-minimal supporting set $L$. For any finite
support $T$, Proposition~\ref{prop:finite-support-intersection} shows that
$L\cap T$ also supports $x$. Its containment in $L$ and the minimality of
$L$ give $L\subseteq L\cap T\subseteq T$. Thus $L$ is least, rather than
merely minimal. Two least supporting sets contain each other and hence
are equal. The final equivalence follows from leastness and enlargement
of supporting bounds.
\end{proof}

Mathlib supplies the well-founded Finset order and
\path{exists_minimal_of_wellFoundedLT}; the nominal step is closure under
binary intersection. The declaration
\path{FinitelySupported.exists_least_support} combines these results.
One could instead minimize the cardinality of a supporting set: intersecting
it with another support cannot decrease its cardinality, so the intersection
equals the selected set. The inclusion-order proof uses the existing order
interface directly. It neither enumerates atoms nor assumes a theorem about
arbitrary intersections.

Let \texttt{hx} be a proof of \texttt{FinitelySupported A x}.
The noncomputable operation \texttt{hx.support} chooses the uniquely
characterized set. Its public laws
state that this set supports $x$ and is contained in every finite support;
clients need not unfold the choice. No public decidable-equality argument is
required: classical equality is local to the intersection proof. Choosing a
semantic support is not an algorithm for computing it.

Two certificates of finite supportedness, even if proved using different
finite bounds, yield the same least set. Lean's proof irrelevance identifies
the certificate arguments, and the mathematical uniqueness result identifies
any independently constructed least support. Neither explanation identifies
supports belonging to different selected actions.

For a nominal carrier, the convenient operation \texttt{support A x} uses
the class's finite-supportedness proof as its elementwise certificate.
It makes no second choice. The agreement law \texttt{support\_eq} equates
this value with \texttt{hx.support} for any supplied certificate of $x$.
Thus two nominality certificates for the same action also yield equal
supports. The atom type is an explicit argument to the carrier operation
and its laws, allowing clients to select the intended action when several
atom types occur in a context.

The carrier interface makes support available uniformly, while the
elementwise interface retains the weaker hypothesis needed for a particular
value. For instance, under the ordinary pointwise action on functions
$\mathbb N\to\mathbb N$, a constant function is supported by its value's
singleton. The identity function has no finite support: given any proposed
bound, swap two distinct atoms outside it and evaluate at one endpoint.
Consequently, this function carrier is not nominal even though its supported
elements admit the least-support construction. No function-space action is
changed to obtain this distinction.

The atom hypothesis cannot be replaced by carrier-wide nominality. Under the
canonical action on $\mathbb B$, every atom has a singleton supporting bound.
Yet the element $\mathrm{false}$ of
Example~\ref{ex:finite-atom-intersection} has no least finite support: a least
bound would lie in both singleton supports, hence be empty, whereas the empty
set does not support it. Finite supportedness and least-support existence
therefore have different atom assumptions.

\begin{proposition}[Laws of least support]
\label{prop:least-support-laws}
Assume $A$ is infinite, $x\in X$ is finitely supported, and $\pi\in\Perm(A)$.
Then
\begin{align*}
  \operatorname{supp}_A(\pi\act x)&=\pi[\operatorname{supp}_A(x)],\\
  \operatorname{supp}_A(x)=\varnothing
    &\quad\Longleftrightarrow\quad
      \forall\sigma\in\Perm(A),\ \sigma\act x=x.
\end{align*}
If $f:X\to Y$ is equivariant, then $f(x)$ is finitely supported and
\[
  \operatorname{supp}_A(f(x))\subseteq\operatorname{supp}_A(x).
\]
No nominality assumption on either whole carrier is needed.
\end{proposition}

\begin{proof}
Transport of sufficient bounds shows that
$\pi[\operatorname{supp}_A(x)]$ supports $\pi\act x$, giving one inclusion
by leastness. For the reverse inclusion, transport the least support of
$\pi\act x$ back by $\pi^{-1}$. It supports $x$, so contains
$\operatorname{supp}_A(x)$; applying $\pi$ yields the other inclusion.
Thus both directions reuse Proposition~\ref{prop:support-transport}.
For the empty characterization, the empty set is a supporting bound exactly
for elements fixed by every permutation, and it is then the least bound.
Finally, equivariance makes every support of $x$ a support of $f(x)$.
Apply this fact to $\operatorname{supp}_A(x)$ and use leastness of the image.
\end{proof}

The elementwise Lean laws take the original certificate \texttt{hx} and
derive those for the renamed element and equivariant image using
\texttt{hx.smul} and \texttt{hx.map}. A separately supplied certificate
gives the same support. The finite-set image equation uses the existing
\texttt{Pointwise} scope and requires \texttt{DecidableEq A}; the empty
characterization and image inclusion need no equality instance. None of
the laws assumes that permutation actions commute.

The carrier versions use the same elementwise laws. In particular, writing
both sides of the image bound with \texttt{support A} requires nominality of
both carriers. For an arbitrary acted-on codomain, the elementwise image
certificate still suffices. Infinitude is needed for both choices of interface;
decidable equality remains confined to the finite-set image equation.

The image inclusion is generally strict. An atom moved by a swap has
nonempty least support, while its image under a constant map into explicit
discrete data is fixed by every permutation and has empty support. Thus
equivariance preserves supporting bounds but need not preserve dependence
on every atom in the original least bound.

Least support is also distinct from strong support. If $\pi\act x=x$,
transport gives
$\pi[\operatorname{supp}_A(x)]=\operatorname{supp}_A(x)$.
This is setwise preservation, which allows atoms inside the support to be
permuted. The supporting property says that fixing those atoms pointwise
fixes $x$; it does not give the converse.
Example~\ref{ex:unordered-pair-support} makes this distinction concrete using
the exact least support of a finite atom set.

\section{Predicates, support and binding}
\label{sec:perspectives}

The distinction between a carrier and its action becomes particularly useful
when the carrier represents logical information. Giving truth values the
trivial action still leaves a nontrivial action on predicates. For a
permutation set $X$, conjugation on $P : X \to \mathrm{Prop}$ gives
\[
  (\pi \act P)(x) \;\longleftrightarrow\; P(\pi^{-1}\act x).
\]
For example, the atom predicate $P_a(x) \equiv (x=a)$ is sent to
$P_{\pi a}$. The singleton $\{a\}$ supports $P_a$, although exchanging $a$
with another atom changes the predicate. Pitts' nominal powerset captures
precisely this distinction: its elements are supported subsets, while
equivariant subsets have empty support~\cite[Section~2.5]{pitts2013}.
There are consequently two different roles for a predicate in reasoning
about binders. As a semantic object, it may carry a finite-support
hypothesis; as the motive of an induction theorem, it may be arbitrary.
Pitts' alpha-structural induction theorem uses supported predicate
families~\cite[Theorem~8.21]{pitts2013}. Nominal2's strong induction principle
instead generalizes each recursive hypothesis over a finitely supported
context, without requiring finite support of the motive
itself~\cite[Section~7, equation~(7.1)]{urbanKaliszyk2012}. These are different
theorem interfaces: the context specifies what binders must avoid, and the
generalized hypotheses provide the strength needed to change that context.
A calculus of supported predicates therefore need not restrict the ordinary
propositions available to an induction proof.

The treatment of support witnesses gives another useful comparison.
The constructive Rocq account of Paranhos and Ventura equips nominal
carriers with a function supplying a finite support and uses setoid equality
with explicit respectfulness laws~\cite[Section~3]{paranhosVentura2025}.
In the present formulation, finite supportedness is an existential
proposition about an element of a selected action. A proof can introduce a
convenient bound, enlarge it, transport it or combine it with another bound,
without making that choice part of the element's identity. This distinction
matters even in classical mathematics: a supplied support may contain atoms
on which the element does not depend. Over infinite atoms, the intersection
theorem turns an inclusion-minimal bound into the unique least support
(Proposition~\ref{prop:least-support}), as in the classical theory of
support~\cite[Section~2.1]{pitts2013}. This construction applies to an
individually supported element; nominality of the whole carrier only supplies
such certificates uniformly. The elementary support calculus itself requires
no operation selecting a least bound. Lean's ordinary
equality and classical foundations make extensional equality and quotient
constructions available; the Rocq formulation instead exposes the chosen
equivalence relation in its action laws. Support arguments must respect the
equality used for the carrier in either setting.

Nominal2 also illustrates the additional mathematical content needed to
connect such a semantic foundation to a language specification. Urban and
Kaliszyk construct raw syntax, define alpha-equivalence, and pass to
quotients, deriving support equations and strong induction through generated
proofs~\cite[Sections~4--7]{urbanKaliszyk2012}. This passage depends on
properties of the particular binding specification: alpha-equivalence must
be an equivariant equivalence relation, and operations descending to the
quotient must respect it. Identifying semantic support with free atoms is a
further theorem about that syntax. The selected-action support laws studied
here supply reusable premises for these arguments, independently of the
constructors chosen for a language. They do not determine which occurrences
a constructor binds. In particular, an action on terms and finite support
alone do not justify a fresh-binder induction rule: the rule must explain
how changing a representative preserves the term and supplies the required
recursive cases. Separating these obligations gives a precise connection
between permutation semantics and the familiar practice of choosing bound
names fresh for the parameters of a proof.


\begin{thebibliography}{9}
\bibitem{pitts2013}
Andrew M. Pitts.
\newblock \emph{Nominal Sets: Names and Symmetry in Computer Science}.
\newblock Cambridge Tracts in Theoretical Computer Science, volume 57.
\newblock Cambridge University Press, 2013.
\bibitem{urbanKaliszyk2012}
Christian Urban and Cezary Kaliszyk.
\newblock General bindings and alpha-equivalence in Nominal Isabelle.
\newblock \emph{Logical Methods in Computer Science}, 8(2:14):1--35, 2012.
\newblock \url{https://doi.org/10.2168/LMCS-8(2:14)2012}.

\bibitem{paranhosVentura2025}
Fabr{\'i}cio Sanches Paranhos and Daniel Ventura.
\newblock Nominal sets in Rocq.
\newblock \emph{Electronic Proceedings in Theoretical Computer Science},
430:55--68, 2025.
\newblock \url{https://doi.org/10.4204/EPTCS.430.5}.
\newblock Also available as arXiv:2509.25883v1.
\end{thebibliography}

\end{document}
